\section{Making the comparison possible}
\label{sec:infrastructure}

The defects in Sections~\ref{sec:instrument} and~\ref{sec:implementation} share a
shape: each is a place where seven independent implementations were free to
differ, and did. Our response is to remove the freedom where it is not part of
the research question, and to make the remaining differences fail loudly.

\subsection{A written specification}

We state the experiment once, as a document rather than as folklore: model
source, optimisation, data pipeline, backend, measurement and replication.
Anything the specification does not pin is a documented degree of freedom, and
anything it pins is checkable.

\subsection{One measurement bridge and one run contract}

The five per-stack CodeCarbon bridges are replaced by one. Every ecosystem drives
it through a line protocol (\code{START}, \code{STOP}, \code{METRIC},
\code{EXIT}), all of which are synchronous, and reads its parameters from a
shared environment contract: output directory, ecosystem, model, dataset,
repetition, seed, epoch count, data and model roots, and the interpreter to
measure with.

Environment variables rather than command-line flags: adding argument parsing to
a C++ binary, a Maven target, an \code{Rscript} and a Rust binary is four
different pieces of plumbing, and four places for the stacks to drift apart
again.

Two things the contract makes impossible rather than merely discouraged. The
interpreter is named explicitly, so no stack can silently be measured by whatever
\code{python3} the shell resolves. And the harness refuses to start when the
framework cannot see the accelerator, so a CPU fallback aborts the run instead of
being recorded as an ecosystem property.

\subsection{One model where the backend allows it}

For the stacks that share LibTorch, the model is exported once as TorchScript
from \code{torchvision} and loaded by each binding. Verified end to end: all six
modules load into Rust and take an optimiser step through the variable store,
with parameter counts matching the Python export exactly
(\num{11227812} for ResNet-18/CIFAR-100), and all six load and forward in R.

R is a documented exception. In \code{torch} 0.17.0 a script module's
\code{\$train()} and \code{\$eval()} raise \code{unused argument}, and the
underlying handle is not reachable through the object's documented fields, so a
loaded module cannot be switched between training and evaluation. Since the
shared modules are exported in training mode, this stack would evaluate with
batch normalisation using batch statistics --- exactly the TensorFlow defect of
Section~\ref{sec:implementation}. It therefore builds its own model, and the
architecture parity that holds for Python, C++ and Rust does not hold for R.
That is a property of the binding, and it belongs in the threats to validity
rather than being papered over.

\subsection{An automated conformance checker}

The specification is enforced by a checker that reads the source the campaign
actually runs: 56 checks covering learning rate, batch size, loader threads, input
scaling, evaluation batching, shared-module use, backend version alignment,
bridge use, tracker synchronisation, device assertion and the absence of
hard-coded paths.

The checker is the piece we would keep if we could keep only one. It turns
``the stacks should be configured the same'' from an intention into a claim that
fails when it stops being true.

\subsection{Every ecosystem verified on the accelerator}

All seven in-scope ecosystems now run on the GPU under the shared contract and
write per-epoch accuracy. Getting there required per-stack fixes that no
documentation mentions: the linker must be told to keep \code{torch\_cuda} or the
Rust binary silently runs on the CPU; \code{CUDA::nvToolsExt} must be defined
before \code{find\_package(Torch)} because NVTX left the CUDA toolkit in CUDA 12;
the C++ embedded interpreter resolves \code{site-packages} from whichever
libpython it was linked against rather than from the campaign environment; R's
\code{liblantern.so} links against \code{libcudart.so.12} while the R bundle ships
it under a hashed name; and DL4J 1.0.0-M2.1 links against the CUDA 11.6 runtime
and is the last release of that API, so on a CUDA 12 host the backend fails with
\code{libcudart.so.11.0: cannot open shared object file} and ND4J reports only
``no backend on your classpath''.

\begin{table}[t]
\centering
\caption{Test accuracy after one epoch on CIFAR-100, identical settings and seed.
The first campaign could not perform this check: no ecosystem wrote its accuracy
to disk.}
\label{tab:convergence}
\small
\begin{tabular}{lrr}
\toprule
Ecosystem & Test accuracy & Train loss \\
\midrule
Python/PyTorch & \SI{19.84}{\percent} & 3.828 \\
C++/LibTorch   & \SI{20.01}{\percent} & 3.822 \\
Rust/tch       & \SI{20.62}{\percent} & 3.798 \\
R/torch        & \SI{20.79}{\percent} & 3.687 \\
Java/DL4J      & \SI{17.14}{\percent} & 3.784 \\
\bottomrule
\end{tabular}
\end{table}

Table~\ref{tab:convergence} is the check that matters. Five independent
implementations agree within \SI{3.6}{\percent}, and the four sharing the
LibTorch backend agree within one point. Java sits lower, which is consistent
with it building its own graph rather than loading the shared module --- and
that is now something a reader can see rather than something nobody recorded.
